\section{Arrays and Vectors}

\subsection{Section Overview}
In this section, we'll explore two fundamental ways to store collections of data in C++: arrays and vectors. We'll start with C-style arrays, which are a fixed-size sequence of elements. Then we'll move on to \texttt{std::vector}, a dynamic and more flexible container from the C++ Standard Library.

\subsection{What is an Array?}
An \textbf{array} is a data structure that stores a fixed-size sequential collection of elements of the same type. Think of it as a row of boxes, where each box holds a value, and the boxes are numbered consecutively starting from 0.

\subsection{Declaring and Initializing Arrays}
You declare an array by specifying the type of its elements, its name, and its size in square brackets.

\begin{lstlisting}
	// Declaring an array of 5 integers
	int scores[5];
\end{lstlisting}

You can initialize an array at declaration using a braced list of initializers.

\begin{lstlisting}
	// Declaring and initializing an array
	int scores[5] {100, 95, 88, 76, 92};

	// The compiler can deduce the size
	int temperatures[] {72, 75, 68, 65}; // size is 4

	// You can also initialize with fewer elements
	double salaries[10] {80000.0, 95000.0}; // remaining elements are 0.0
\end{lstlisting}

\subsection{Accessing and Modifying Array Elements}
You access elements in an array using the subscript operator \texttt{[]}. Array indices in C++ start at 0.

\begin{lstlisting}
	int scores[] {100, 95, 88};

	// Accessing elements
	std::cout << "First score: " << scores[0] << std::endl;  // 100
	std::cout << "Second score: " << scores[1] << std::endl; // 95
	std::cout << "Third score: " << scores[2] << std::endl;  // 88

	// Modifying elements
	scores[1] = 98;
	std::cout << "New second score: " << scores[1] << std::endl; // 98
\end{lstlisting}

\begin{remark}
	C++ does not perform bounds checking for arrays. Accessing an array beyond its boundaries (e.g., \texttt{scores[5]} in an array of size 5) leads to undefined behavior. This is a common source of bugs.
\end{remark}

\subsection{Coding Exercise 5: Declaring, Initializing and Accessing an Array}
Write a C++ program that declares an array of 3 \texttt{double} values named \texttt{heights}. Initialize it with the values \texttt{1.75}, \texttt{1.82}, and \texttt{1.69}. Then, print the second element of the array to the console.

\begin{lstlisting}
	#include <iostream>

	int main() {
		double heights[3] {1.75, 1.82, 1.69};
		std::cout << "The second height is: " << heights[1] << std::endl;
		return 0;
	}
\end{lstlisting}

\subsection{Multidimensional Arrays}
C++ supports multidimensional arrays, which are arrays of arrays. A 2D array is like a grid or a table.

\begin{lstlisting}
	// Declaring a 2D array (3 rows, 4 columns)
	int movie_ratings [3][4] {
		{0, 4, 3, 5},
		{1, 3, 4, 4},
		{2, 4, 4, 5}
	};

	// Accessing an element (row 1, column 2)
	std::cout << "Rating: " << movie_ratings[1][2] << std::endl; // 4
\end{lstlisting}

\subsection{Declaring and Initializing Vectors}
The C++ Standard Library provides \texttt{std::vector}, a dynamic array that can grow or shrink in size. It's generally safer and more flexible than C-style arrays. To use \texttt{std::vector}, you must include the \texttt{<vector>} header.

\begin{lstlisting}
	#include <vector>

	// Declaring a vector of integers
	std::vector<int> scores;

	// Declaring and initializing a vector
	std::vector<int> temperatures {72, 75, 68, 65};

	// Initialize with a certain size and value
	std::vector<double> salaries(10, 80000.0); // 10 doubles, all 80000.0
\end{lstlisting}

\subsection{Accessing and Modifying Vector Elements}
You can access vector elements using the subscript operator \texttt{[]} just like arrays. \texttt{std::vector} also provides the \texttt{.at()} method, which performs bounds checking.

\begin{lstlisting}
	#include <iostream>
	#include <vector>

	int main() {
		std::vector<int> scores {100, 95, 88};

		// Accessing with []
		std::cout << "First score: " << scores[0] << std::endl;

		// Accessing with .at()
		std::cout << "Second score: " << scores.at(1) << std::endl;

		// Modifying elements
		scores.at(1) = 98;
		std::cout << "New second score: " << scores.at(1) << std::endl;

		// Adding elements
		scores.push_back(80); // adds 80 to the end

		std::cout << "Size of vector: " << scores.size() << std::endl; // 4

		return 0;
	}
\end{lstlisting}

\begin{remark}
	The \texttt{.at()} method will throw an exception if you try to access an element out of bounds, which can help you catch errors. The \texttt{[]} operator does not provide this safety.
\end{remark}

\subsection{Coding Exercise 6: Declaring, Initializing and Accessing Vectors}
Write a C++ program that declares a vector of \texttt{std::string} named \texttt{planets}. Initialize it with ``Mercury'', ``Venus'', and ``Earth''. Then, add ``Mars'' to the end of the vector and print the size of the vector.

\begin{lstlisting}
	#include <iostream>
	#include <vector>
	#include <string>

	int main() {
		std::vector<std::string> planets {"Mercury", "Venus", "Earth"};
		planets.push_back("Mars");
		std::cout << "The number of planets is: "
			<< planets.size() << std::endl;
		return 0;
	}
\end{lstlisting}

\subsection{Section Challenge}
Write a program that declares two empty vectors of integers named \texttt{vector1} and \texttt{vector2}.
\begin{enumerate}
	\item Add 10 and 20 to \texttt{vector1} dynamically using \texttt{push\_back}.
	\item Display the elements in \texttt{vector1} using the \texttt{at()} method as well as its size using the \texttt{size()} method.
	\item Add 100 and 200 to \texttt{vector2} dynamically using \texttt{push\_back}.
	\item Display the elements in \texttt{vector2} using the \texttt{at()} method as well as its size using the \texttt{size()} method.
	\item Declare an empty 2D vector called \texttt{vector\_2d}.
	\item Add \texttt{vector1} and \texttt{vector2} to \texttt{vector\_2d} dynamically using \texttt{push\_back}.
	\item Display the elements in \texttt{vector\_2d} using the \texttt{at()} method.
	\item Change \texttt{vector1.at(0)} to 1000.
	\item Display the elements in \texttt{vector\_2d} again.
\end{enumerate}

\subsection{Section Challenge --- Solution}
\begin{lstlisting}
	#include <iostream>
	#include <vector>

	int main() {
		std::vector<int> vector1;
		std::vector<int> vector2;

		vector1.push_back(10);
		vector1.push_back(20);

		std::cout << "vector1 elements: "
			<< vector1.at(0) << " " << vector1.at(1) << std::endl;
		std::cout << "vector1 size: " << vector1.size() << std::endl;

		vector2.push_back(100);
		vector2.push_back(200);

		std::cout << "vector2 elements: "
			<< vector2.at(0) << " " << vector2.at(1) << std::endl;
		std::cout << "vector2 size: " << vector2.size() << std::endl;

		std::vector<std::vector<int>> vector_2d;
		vector_2d.push_back(vector1);
		vector_2d.push_back(vector2);

		std::cout << "vector_2d elements:" << std::endl;
		std::cout << vector_2d.at(0).at(0) << " "
			<< vector_2d.at(0).at(1) << std::endl;
		std::cout << vector_2d.at(1).at(0) << " "
			<< vector_2d.at(1).at(1) << std::endl;

		vector1.at(0) = 1000;

		std::cout << "vector_2d elements after change:" << std::endl;
		std::cout << vector_2d.at(0).at(0) << " "
			<< vector_2d.at(0).at(1) << std::endl;
		std::cout << vector_2d.at(1).at(0) << " "
			<< vector_2d.at(1).at(1) << std::endl;

		std::cout << "vector1 elements after change: "
			<< vector1.at(0) << " " << vector1.at(1) << std::endl;

		return 0;
	}
\end{lstlisting}

\begin{remark}
	Notice that when we changed \texttt{vector1}, the \texttt{vector\_2d} did NOT change. This is because \texttt{vector\_2d} contains copies of \texttt{vector1} and \texttt{vector2}.
\end{remark}

\subsection{Quiz 4: Section 07 Quiz}
\begin{enumerate}
	\item What is the index of the first element in a C++ array?
	\begin{enumerate}
		\item[a)] 1
		\item[b)] 0
		\item[c)] $-1$
		\item[d)] It depends on the declaration.
	\end{enumerate}

	\item What is a major advantage of \texttt{std::vector} over C-style arrays?
	\begin{enumerate}
		\item[a)] They are faster.
		\item[b)] They can dynamically resize.
		\item[c)] They use less memory.
		\item[d)] They are part of the C language.
	\end{enumerate}

	\item Which method should you use to access a \texttt{std::vector} element if you want bounds checking?
	\begin{enumerate}
		\item[a)] \texttt{[]}
		\item[b)] \texttt{get()}
		\item[c)] \texttt{at()}
		\item[d)] \texttt{element()}
	\end{enumerate}
\end{enumerate}

\textbf{Answers:} 1.\ (b), 2.\ (b), 3.\ (c)
